\section{Methodology}

\ProjectFigure
  {figures/phase_03/fig_p03_01_experiment_pipeline.pdf}
  {0.96\columnwidth}
  {Pipeline diagram showing preprocessing, optional training augmentation,
  three model families, and clean/20/10/0 dB evaluation -- pending}
  {Overview of the shared training and controlled robustness-evaluation
  pipeline. Augmentation is applied only to training data.}
  {fig:experiment-pipeline}

\subsection{Audio Preprocessing}

All recordings were converted to mono, resampled to 22,050 Hz, and normalized
to four seconds. Short recordings were zero-padded. Training used seeded random
crops, whereas validation and test processing used deterministic center crops.
The common model input was a per-example standardized Log-Mel spectrogram with
shape $1 \times 64 \times 173$.

\begin{table}[t]
  \caption{Shared Preprocessing Configuration}
  \label{tab:preprocessing}
  \centering
  \begin{tabular}{ll}
    \toprule
    Setting & Value \\
    \midrule
    Sample rate & 22,050 Hz \\
    Clip duration & 4.0 s \\
    Input channels & Mono \\
    FFT/window length & 1,024 \\
    Hop length & 512 \\
    Mel bands & 64 \\
    Dynamic range & 80 dB \\
    Input shape & $1 \times 64 \times 173$ \\
    Normalization & Per-example standardization \\
    \bottomrule
  \end{tabular}
\end{table}


\subsection{Training-Time Augmentation}

The baseline condition disabled robustness-oriented augmentation. The augmented
condition independently applied time shift, random gain, background-noise
mixing, frequency masking, and time masking with probability 0.5. Background
noise used a training-only MS-SNSD noise bank and a uniformly selected SNR in
the configured 0--20 dB range.

\subsection{Model Architectures}

\begin{table}[t]
  \caption{Model Configurations}
  \label{tab:model-parameters}
  \centering
  \begin{tabular}{lrr}
    \toprule
    Architecture & Batch size & Parameters \\
    \midrule
    CNN & 16 & 94,186 \\
    CRNN--BiGRU & 16 & 293,610 \\
    ResNet18 & 8 & 11,175,370 \\
    \bottomrule
  \end{tabular}
\end{table}


Every architecture accepted the same one-channel Log-Mel representation and
returned ten class logits. The ResNet18 model was trained from scratch rather
than initialized with ImageNet weights.

\subsection{Controlled Noise Corruption}

Final evaluation used clean audio and noisy variants at 20, 10, and 0 dB. Noise
was scaled using mean-square signal and noise power:

\begin{equation}
  \mathrm{SNR}_{\mathrm{dB}}
  = 10 \log_{10}\left(\frac{P_{\mathrm{signal}}}{P_{\mathrm{noise}}}\right).
\end{equation}

The corruption seed and sample identifier deterministically selected the noise
file and segment. Therefore, every model received the same corruption for a
given test sample and condition. Training noise and evaluation noise came from
separate MS-SNSD directories.

\ProjectFigure
  {figures/phase_02/fig_p02_08_snr_waveform_comparison.png}
  {\columnwidth}
  {Controlled SNR waveform comparison}
  {Example waveform corruption under the configured clean, 20 dB, 10 dB, and
  0 dB conditions.}
  {fig:snr-waveform-comparison}

\subsection{Metrics}

The study reports accuracy, macro precision, macro recall, macro F1, per-class
precision/recall/F1, and confusion matrices. Robustness summaries include the
clean-to-noisy drop, relative retention, ordinary least-squares slope over the
finite SNR points, and normalized trapezoidal area under the 0--20 dB SNR curve.
